\begin{textsample}{NEG Dim 1 – human – Score -7.00 – mg/mg\_ef\_linguagens\_ed\_fis\_12\_p1.txt}  \label{ex:f1_neg_004}
Também estruturadas a partir de culturas diversas, as lutas compõem nosso repertório de práticas \textbf{corporais} com \textbf{sentidos}, \textbf{significados} e tradições vinculados às várias \textbf{formas} de resistência, disputas, códigos, culturas e comportamentos.
São compostas por \textbf{movimentos} nos “ [... ] quais os participantes empregam técnicas, táticas e estratégias específicas para imobilizar, desequilibrar, atingir ou excluir o oponente de um \textbf{determinado} espaço, combinando ações de ataque e defesa [... ] ” ( BNCC, 2017 ).
Possuem também um caráter de valores próprios, atribuído a cada uma de suas manifestações que deve ser contemplado e \textbf{respeitado}, mas também refletido e recriado.
Ao trabalhar com as habilidades dessa unidade temática é importante estar comprometido com a origem, história e elementos constituintes, garantindo que os estudantes possam compreender essa manifestação \textbf{corporal} de \textbf{forma} profunda e crítica, estando apto a agir para a mediação de disputas e conflitos de \textbf{forma} ética, honrosa e pacífica.

% matched lemmas: corporal, determinar, forma, movimento, respeitar, sentido, significado
\end{textsample}
